
\section{Hofer 几何}\label{sec:9-6}

Hofer 在 \srcite{177} 中定义了欧氏空间之紧支撑\BookRecordIndex{T0362}辛微分同胚群上的一个度量，并通过研究
相关哈密顿流的周期轨道发展了它的性质。Lalonde--McDuff \srcite{223,224} 采用了更具
几何味道的方法，把 \BookRecordIndex{T0027}Hofer 度量与\BookRecordIndex{T0156}哈密顿纤维化的若干问题联系起来。本节将解释
Polterovich 对这些思想的一项优美推广。我们的讨论会保持相当集中的范围；关于这个主题
的一般背景，读者可查阅 Polterovich \srcite{330,328} 或
\srcite{277, Chapter 12}。以下始终假设 $(M,\omega)$ 是紧\BookRecordIndex{T0358}辛流形。讨论中最重要的部分
将考察单调情形，不过下面说明的许多结果在更一般的情形仍成立。

设
\[
  \mathbb R\longrightarrow\operatorname{Symp}(M,\omega),
  \qquad t\longmapsto\psi_t=\psi_{t+1}
\]
是\BookRecordIndex{T0161}哈密顿辛微分同胚的环路。则存在光滑哈密顿函数
\[
  \mathbb R\longrightarrow C^\infty(M),
  \qquad t\longmapsto H_t,
\]
使每个 $t\in\mathbb R$ 都满足
\[
  \frac d{dt}\psi_t=X_{H_t}\circ\psi_t,
  \qquad H_t=H_{t+1}.
\]
为以后使用，最好把哈密顿函数归一化为平均值全为零：
\[
  \int_M H_t\omega^n=0.
\]
环路 $t\mapsto\psi_t$ 的 \emph{Hofer 长度}定义为
\[
  \ell(\{\psi_t\}):=\int_0^1\|H_t\|\,dt,
  \qquad
  \|H_t\|:=\max_MH_t-\min_MH_t.
\]
注意，这个长度不依赖上述归一化。

\begin{example}\label{ex:9-6-1}
设 $\omega_{\mathrm{FS}}\in\Omega^2(\mathbb{CP}^n)$ 是体积为 $\pi^n/n!$ 的
Fubini--Study 形式。考察环路
\[
  \mathbb R/\mathbb Z\longrightarrow
  \operatorname{Symp}(\mathbb{CP}^n,\omega_{\mathrm{FS}}),
  \qquad t\longmapsto\psi_t^k,
\]
其中
\[
  \psi_t^k([z_0:\cdots:z_n])
  :=[z_0:e^{2\pi it}z_1:\cdots:e^{2\pi it}z_k:z_{k+1}:\cdots:z_n].
\]
该环路由与时间无关的哈密顿函数
\begin{equation}\label{eq:9-6-1}
  H^k([z_0:\cdots:z_n])
  :=\pi\left(
    \frac{k}{n+1}
    -\frac{|z_1|^2+\cdots+|z_k|^2}
      {|z_0|^2+\cdots+|z_n|^2}
  \right)
\end{equation}
生成。由于 $\|H^k\|=\pi$，可得
\[
  \ell(\{\psi_t^k\})=\pi.
\]
Lalonde--McDuff 在 \srcite{224} 中证明，当 $n=k=1$ 时，环路
$t\mapsto\psi_t$ 不能通过哈密顿变形缩短。他们的论证使用嵌入球和以 $S^2$ 为底的
\BookRecordIndex{T0156}哈密顿纤维化之\BookRecordIndex{T0120}非挤压定理。本节的目的是解释 Polterovich 给出的另一种证明：它使用例
~\ref{ex:8-6-8} 的 \BookRecordIndex{T0047}Seidel 表示，并且对所有 $1\le k\le n$ 都成立。我们还将说明环路的
另一些几何定义的不变量；它们最终都与该环路的 Hofer 长度相等。
\end{example}

Polterovich 在 \srcite{325,326,327} 中考察了 $M$ 上\BookRecordIndex{T0161}哈密顿辛微分同胚环路
$\psi=\{\psi_t\}_{t\in\mathbb R/\mathbb Z}$ 所关联的三个不变量。第一个是所有与
$\psi$ \BookRecordIndex{T0158}哈密顿同痕的哈密顿环路之最小 Hofer 长度：
\[
  \lambda(\psi):=\inf_{\psi'\sim\psi}\ell(\psi').
\]
第二个是在由式~\eqref{eq:8-2-6} 定义的\BookRecordIndex{T0156}哈密顿纤维化
$\widetilde M_\psi$ 上，所有\BookRecordIndex{T0224}联络形式之曲率 $R_{\widetilde\omega}$ 的 \BookRecordIndex{T0028}Hofer 范数下确界：
\[
  \kappa(\psi):=
  \inf_{\widetilde\omega\in\mathcal T_0(\widetilde M_\psi)}
  \|R_{\widetilde\omega}\|.
\]
这里 $\mathcal T_0(\widetilde M_\psi)$ 表示注~\ref{rem:8-2-11} 中所有\BookRecordIndex{T0224}联络形式
$\widetilde\omega\in\Omega^2(\widetilde M_\psi)$ 所成空间。第三个是\emph{\BookRecordIndex{T0122}非辛区间的
长度}：
\[
  \varepsilon(\psi):=
  \inf_{\widetilde\omega_c^{n+1}>0}\int_{S^2}c\sigma
  -\sup_{\widetilde\omega_c^{n+1}<0}\int_{S^2}c\sigma.
\]
这里第一个下确界取遍所有\BookRecordIndex{T0224}联络形式
\[
  \widetilde\omega_c
  =\widetilde\omega+\pi^*(c\sigma)
  \in\Omega^2(\widetilde M_\psi),
\]
其中 $\widetilde\omega\in\mathcal T_0(\widetilde M_\psi)$、
$c\in\Omega^0(S^2)$，且 $\widetilde\omega_c^{n+1}>0$（等价地，
$c\circ\pi>R_{\widetilde\omega}$）。上确界作同样理解。

\begin{remark}\label{rem:9-6-2}
这个\BookRecordIndex{T0122}非辛区间长度与第~\ref{sec:9-3} 节末定义的以 $S^2$ 为底之辛纤维丛的面积有关。在
当前情形，该面积为
\[
  \operatorname{area}(\widetilde M_\psi,\widetilde\omega_c)
  :=\frac{\operatorname{Vol}(\widetilde M_\psi,\widetilde\omega_c)}
    {\operatorname{Vol}(M,\omega)}
  =\int_{S^2}c\sigma.
\]
所以
$\varepsilon(\psi)=\varepsilon^+(\psi)+\varepsilon^-(\psi)$，其中
$\varepsilon^+(\psi)$ 是 $\widetilde M_\psi$ 上一个\BookRecordIndex{T0224}联络形式的最小面积，而
$\varepsilon^-(\psi):=\varepsilon^+(\psi^{-1})$。上面定义的另两个不变量也都分解为两
部分之和，因为任意归一化函数 $R$ 的总变差 $\|R\|$ 可写成
$\|R\|_++\|R\|_-$，其中
\[
  \|R\|_+:=\max_MR,
  \qquad \|R\|_-:=-\min_MR=\|-R\|_+.
\]
事实上，Polterovich 的结果讨论的是这些“单侧”不变量。更多细节见 Polterovich
\srcite{325,326,327} 与 McDuff \srcite{271}。
\end{remark}

Polterovich 在 \srcite{325,326,327} 中证明这三个不变量全都相等。我们复述他对下述较弱
结论的证明。

\begin{proposition}[Polterovich]\label{prop:9-6-3}
设 $(M,\omega)$ 是紧\BookRecordIndex{T0358}辛流形，
$\psi=\{\psi_t\}_{t\in\mathbb R/\mathbb Z}$ 是 $M$ 的\BookRecordIndex{T0161}哈密顿辛微分同胚环路。则
\[
  \lambda(\psi)\ge\kappa(\psi)\ge\varepsilon(\psi).
\]
\end{proposition}

\begin{proof}
首先注意，存在\BookRecordIndex{T0224}联络形式
$\widetilde\omega\in\mathcal T_0(\widetilde M_\psi)$，使
\[
  \|R_{\widetilde\omega}\|=\ell(\psi).
\]
\sourcefootnote{原书在这里把\BookRecordIndex{T0224}联络形式所属空间印作
$\mathcal T(\widetilde M_\psi)$，但本节只定义并使用
$\mathcal T_0(\widetilde M_\psi)$；紧随其后的公式也正是注~\ref{rem:8-2-11} 中该空间的
\BookRecordIndex{T0224}联络形式，故补上下标 $0$。}
在注~\ref{rem:8-2-11} 的记号中，一个显式公式是
\[
  F_{s,t}^{\pm}:=0,
  \qquad G_{s,t}^+:=-\rho(s)H_t,
  \qquad G_{s,t}^-:=-\rho(s)H_{-t}\circ\psi_{-t},
\]
其中 $\rho:\mathbb R\to[0,1]$ 是光滑\BookRecordIndex{T0188}截断函数，在 $s\le-1$ 时消失，在 $s\ge1$ 时等于
一，并满足 $\rho(s)+\rho(-s)\equiv1$。

其次，若两个哈密顿环路 $\psi,\psi'$ \BookRecordIndex{T0158}哈密顿同痕，则相应\BookRecordIndex{T0156}哈密顿纤维化
$\widetilde M_\psi,\widetilde M_{\psi'}$（逐纤维）同构。为看出这一点，任取\BookRecordIndex{T0158}哈密顿同痕
\[
  \mathbb R\times\mathbb R/\mathbb Z
  \longrightarrow\operatorname{Symp}(M,\omega),
  \qquad(s,t)\longmapsto\psi_{s,t},
\]
使 $s\le-1$ 时 $\psi_{s,t}=\psi_t$，而 $s\ge1$ 时
$\psi_{s,t}=\psi'_t$。定义微分同胚
$\Psi:\widetilde M_\psi\to\widetilde M_{\psi'}$：在正平凡化上令
\[
  \Psi^+(s,t,x^+)
  :=(s,t,\psi_{s,t}\circ\psi_t^{-1}(x^+)),
\]
在负平凡化上令
\[
  \Psi^-(s,t,x^-)
  :=(s,t,(\psi'_{-t})^{-1}\circ\psi_{-s,-t}(x^-)).
\]
对每个
$\widetilde\omega'\in\mathcal T_0(\widetilde M_{\psi'})$，都有
$R_{\Psi^*\widetilde\omega'}=R_{\widetilde\omega'}\circ\Psi$，故
$\|R_{\Psi^*\widetilde\omega'}\|=\|R_{\widetilde\omega'}\|$。上面已经找到满足
$\|R_{\widetilde\omega'}\|=\ell(\psi')$ 的\BookRecordIndex{T0224}联络形式，因此
$\kappa(\psi)\le\lambda(\psi)$。

第三，证明 $\varepsilon(\psi)\le\kappa(\psi)$。固定 $\delta>0$，并选取\BookRecordIndex{T0224}联络形式
$\widetilde\omega\in\mathcal T_0(\widetilde M_\psi)$，使
\[
  \|R_{\widetilde\omega}\|<\kappa(\psi)+\delta.
\]
则存在光滑函数 $c,K:S^2\to\mathbb R$，使
\[
  c(z)<\min_{\pi^{-1}(z)}R_{\widetilde\omega}(z)
  \le\max_{\pi^{-1}(z)}R_{\widetilde\omega}(z)<K(z),
  \qquad
  \int_{S^2}(K-c)\sigma<\kappa(\psi)+\delta.
\]
记
$\widetilde\omega_c:=\widetilde\omega+\pi^*(c\sigma)$ 与
$\widetilde\omega_K:=\widetilde\omega+\pi^*(K\sigma)$。则
$\widetilde\omega_K^{n+1}>0$ 且
$\widetilde\omega_c^{n+1}<0$，因而
\[
  \varepsilon(\psi)
  \le\int_{S^2}(K-c)\sigma
  <\kappa(\psi)+\delta.
\]
该式对每个 $\delta>0$ 都成立，所以
$\varepsilon(\psi)\le\kappa(\psi)$。命题得证。
\end{proof}

\begin{proposition}[Polterovich--Seidel]\label{prop:9-6-4}
设
\[
  \mathbb{CP}^n\lhook\joinrel\longrightarrow
  \widetilde M^k\longrightarrow S^2
\]
是例~\ref{ex:9-6-1} 中环路 $\psi^k$ 所关联的\BookRecordIndex{T0156}哈密顿纤维化。则对
$k=1,\ldots,n$，存在同调类
$\widetilde A^k\in H_2(\widetilde M^k;\mathbb Z)$，使每个\BookRecordIndex{T0224}联络形式
$\widetilde\omega\in\mathcal T_0(\widetilde M^k)$ 都满足
\begin{equation}\label{eq:9-6-2}
  \mathcal S(\psi^k)=\mathcal S_{\widetilde A^k}(\psi^k)=c^{n+1-k},
  \qquad
  \int_{\widetilde A^k}\widetilde\omega
  =-\frac{\pi(n+1-k)}{n+1}.
\end{equation}
这里 $c\in H^2(\mathbb{CP}^n)$ 表示正生成元，而
$\mathcal S(\psi)\in H^*(\mathbb{CP}^n)$ 是例~\ref{ex:8-6-8} 中由式~\eqref{eq:8-6-4}
定义的。
\end{proposition}

\begin{proof}
环路
$\psi_t^k\in\operatorname{Symp}(\mathbb{CP}^n,\omega_{\mathrm{FS}})$ 保持度量，并由与
时间无关的哈密顿函数 $H^k:\mathbb{CP}^n\to\mathbb R$ 生成（见式
~\eqref{eq:9-6-1}）。因此，标准复结构 $J_0$ 在 $\widetilde M^k$ 上定义了\BookRecordIndex{T0306}竖直殆复结构。
在注~\ref{rem:8-2-11} 的记号中，一个\BookRecordIndex{T0224}联络形式
$\widetilde\omega\in\mathcal T_0(\widetilde M^k)$ 由公式
\begin{equation}\label{eq:9-6-3}
  F^\pm:=0,
  \qquad G_{s,t}^\pm:=-\rho(s)H^k
\end{equation}
给出，其中 $\rho:\mathbb R\to[0,1]$ 是光滑\BookRecordIndex{T0188}截断函数，满足：$s\le-1$ 时
$\rho(s)=0$，$s\ge1$ 时 $\rho(s)=1$，并且对所有 $s$，
$\dot\rho(s)\ge0$ 且 $\rho(s)+\rho(-s)=1$。注意
\[
  R_{\widetilde\omega}\sigma=-\dot\rho(s)H^k\,ds\wedge dt.
\]
以 $\widetilde J$ 表示由 $J_0,\widetilde\omega$ 决定的 $\widetilde M^k$ 上\BookRecordIndex{T0088}殆复结构，并以
$\mathcal M(\widetilde A;J_0,\widetilde\omega)$ 表示\BookRecordIndex{T0156}哈密顿纤维化
$\widetilde M^k\to S^2$ 的 $\widetilde J$-全纯截面空间。

现考察临界流形
\[
  Z^k:=\bigl\{[0:z_1:\cdots:z_k:0:\cdots:0]
  \bigm|(z_1,\ldots,z_k)\in\mathbb C^k\setminus\{0\}\bigr\}
  \subset\mathbb{CP}^n.
\]
函数 $H^k$ 在这个子流形上取得极小值
$\pi(k-n-1)/(n+1)$。每个常值映射
$u^\pm(s,t)\equiv x\in Z^k$ 都定义一个\BookRecordIndex{T0315}水平截面
$\widetilde u:S^2\to\widetilde M^k$，因而是 $\widetilde J$-全纯的。以
$\widetilde A^k\in H_2(\widetilde M^k;\mathbb Z)$ 表示这些常值截面
$\widetilde u$ 的同调类。由引理~\ref{lem:8-2-9}，
\[
\begin{aligned}
  \int_{\widetilde A^k}\widetilde\omega
  &=E^{\mathrm{Vert}}(\widetilde u)
    -\int_{S^2}(R_{\widetilde\omega}\circ\widetilde u)\sigma\\
  &=\int_0^1\int_{-\infty}^{\infty}
      \dot\rho(s)H^k(x)\,ds\,dt
   =\frac{\pi(k-n-1)}{n+1}.
\end{aligned}
\]
这说明式~\eqref{eq:9-6-2} 的第二个等式对某个\BookRecordIndex{T0224}联络形式成立，因而对每个
$\widetilde\omega\in\mathcal T_0(\widetilde M^k)$ 都成立。

考察沿取值于 $Z^k$ 的一个水平解的\BookRecordIndex{T0308}竖直切丛，可见它分裂成复线丛的直和：其中 $k-1$
个的 Chern 数为零（它们构成截面沿 $Z^k$ 切方向的\BookRecordIndex{T0117}法丛），另外 $n+1-k$ 个的 Chern 数
为 $-1$。所以竖直第一 Chern 数为
\[
  c_1^{\mathrm{Vert}}(\widetilde A^k)=k-1-n.
\]
此外，每个这类常值截面的\BookRecordIndex{T0346}线性化算子都是满射。下面证明该同调类中不可能有其他
$\widetilde J$-全纯截面。事实上，若
$\widetilde u=(u^+,u^-)\in
\mathcal M(\widetilde A^k;J_0,\widetilde\omega)$，则由引理~\ref{lem:8-2-9}，
\[
\begin{aligned}
  0&\le E^{\mathrm{Vert}}(\widetilde u)\\
   &=\int_{S^2}\widetilde u^*\widetilde\omega
     +\int_{S^2}(R_{\widetilde\omega}\circ\widetilde u)\sigma\\
   &=\frac{\pi(k-n-1)}{n+1}
     -\int_0^1\int_{-\infty}^{\infty}
       \dot\rho(s)H^k(u^\pm(s,t))\,ds\,dt\\
   &\le\frac{\pi(k-n-1)}{n+1}
     -\int_0^1\int_{-\infty}^{\infty}
       \dot\rho(s)\min_{\mathbb{CP}^n}H^k\,ds\,dt
    =0.
\end{aligned}
\]
因此处处取等号，每条表示类 $\widetilde A^k$ 的 $\widetilde J$-全纯曲线都必须水平，故
必为上面已考察的常值截面之一。（其他\BookRecordIndex{T0315}水平截面对应于 $H^k$ 的极大值；它们的 Chern 数
不同，因而不表示类 $\widetilde A^k$。）
\sourcefootnote{原书把括号中“其他\BookRecordIndex{T0315}水平截面”说成对应 $H^k$ 的极小值；但前文已说明
$Z^k$ 正是极小值流形，而其余固定分支是极大值流形。故把 minima 唯一补正为 maxima。}
现在\BookRecordIndex{T0283}求值映射
\[
  \operatorname{ev}_0:
  \mathcal M(\widetilde A^k;J_0,\widetilde\omega)
  \longrightarrow\mathbb{CP}^n
\]
显然是到 $Z^k$ 上的微分同胚。由于
$Z^k\subset\mathbb{CP}^n$ 是复维 $k-1$ 的射影子空间，便有
\[
  \operatorname{GW}^{\widetilde M^k,\{1\}}_{\widetilde A^k,1}(c^{k-1})=1,
  \qquad \mathcal S_{\widetilde A^k}(\psi^k)=c^{n+1-k}.
\]
为完成证明，还须说明
\[
  \widetilde A\ne\widetilde A^k
  \quad\Longrightarrow\quad
  \mathcal S_{\widetilde A}(\psi^k)=0.
\]
注意，若 $\mathcal S_{\widetilde A}(\psi^k)\ne0$，则 $\widetilde A$ 必须满足
\[
  \pi_*\widetilde A=[S^2],
  \qquad -n\le c_1^{\mathrm{Vert}}(\widetilde A)\le0.
\]
当 $1\le k\le n$ 时，唯一这样的类就是 $\widetilde A^k$。命题得证。
\end{proof}

\begin{theorem}[Polterovich]\label{thm:9-6-5}
设 $\psi^k=\{\psi_t^k\}_{t\in\mathbb R/\mathbb Z}$ 是例~\ref{ex:9-6-1} 中讨论的环路。
若 $k\in\{1,\ldots,n\}$，则
\[
  \lambda(\psi^k)=\kappa(\psi^k)=\varepsilon(\psi^k)=\pi.
\]
\end{theorem}

\begin{proof}
由命题~\ref{prop:9-6-3}，
\[
  \varepsilon(\psi^k)\le\kappa(\psi^k)
  \le\lambda(\psi^k)\le\pi.
\]
所以只须证明
$\varepsilon(\psi^k)\ge\pi$。为此，设
$\widetilde\omega\in\mathcal T_0(\widetilde M^k)$ 与
$c\in\Omega^0(S^2)$ 满足
\[
  \widetilde\omega_c^{n+1}>0,
  \qquad \widetilde\omega_c:=\widetilde\omega+c\sigma.
\]
令 $\widetilde A^k\in H_2(\widetilde M^k;\mathbb Z)$ 是命题~\ref{prop:9-6-4} 中的同调
类，并固定 $\mathbf w\in S^2$。由定理~\ref{thm:8-5-3}，正则偶对集合
$\mathcal{JH}_{\mathrm{reg}}(\widetilde M^k,\pi,\widetilde\omega;\mathbf w)$ 在
$\mathcal J^{\mathrm{Vert}}\times\mathcal H$ 中关于 $C^\infty$ 拓扑稠密。因此，对
$\widetilde\omega$ 作任意小的\BookRecordIndex{T0402}正合扰动后，可假设存在\BookRecordIndex{T0306}竖直殆复结构
$J\in\mathcal J^{\mathrm{Vert}}(\widetilde M^k,\pi,\widetilde\omega)$，使偶对 $(J,0)$
对 $\mathbf w$ 正则（定义~\ref{def:8-5-2} 的意义下），并且条件
$\widetilde\omega_c^{n+1}>0$ 仍然满足。记
$\mathcal M(\widetilde A^k;J,\widetilde\omega)$ 为 $\widetilde M^k$ 中表示类
$\widetilde A^k$ 的 $\widetilde J$-全纯截面之\BookRecordIndex{T0246}模空间，其中 $\widetilde J$ 是由
$J,\widetilde\omega$ 诱导的 $\widetilde M^k$ 上\BookRecordIndex{T0088}殆复结构（见引理~\ref{lem:8-2-8}）。
\sourcefootnote{原书在此处\BookRecordIndex{T0283}求值映射的定义域把\BookRecordIndex{T0246}模空间类漏写了上标 $k$，印作
$\mathcal M(\widetilde A;J,\widetilde\omega)$；本段始终固定命题~\ref{prop:9-6-4} 的类
$\widetilde A^k$，故补全。}
于是由定理~\ref{thm:8-5-1}，\BookRecordIndex{T0283}求值映射
\[
  \widetilde{\operatorname{ev}}_{\mathbf w}:
  \mathcal M(\widetilde A^k;J,\widetilde\omega)
  \longrightarrow\mathbb{CP}^n,
  \qquad
  \widetilde{\operatorname{ev}}_{\mathbf w}(\widetilde u)
  :=\widetilde\iota_{\mathbf w}^{-1}(\widetilde u(\mathbf w))
\]
是\BookRecordIndex{T0337}伪循环；由命题~\ref{prop:9-6-4}，它在 $\mathbb{CP}^n$ 中表示非零同调类。因此
$\mathcal M(\widetilde A^k;J,\widetilde\omega)\ne\varnothing$。由于
$\widetilde\omega_c$ 是 $\widetilde M^k$ 上的\BookRecordIndex{T0363}辛形式，这蕴含
\[
  0\le\langle[\widetilde\omega_c],\widetilde A^k\rangle
  =\int_{S^2}c\sigma-\frac{\pi(n+1-k)}{n+1},
\]
其中第二个等式来自式~\eqref{eq:9-6-2}。因此已证明
\[
  \inf_{\substack{
    \widetilde\omega\in\mathcal T_0(\widetilde M^k),\ c\in\Omega^0(S^2)\\
    \widetilde\omega_c^{n+1}>0}}
  \int_{S^2}c\sigma
  \ge\frac{\pi(n+1-k)}{n+1}.
\]
由于环路 $t\mapsto\psi_{-t}^k$ \BookRecordIndex{T0158}哈密顿同痕于
$t\mapsto\psi_t^{n+1-k}$，可得
\[
\begin{aligned}
  -\sup_{\substack{
    \widetilde\omega\in\mathcal T_0(\widetilde M^k),\ c\in\Omega^0(S^2)\\
    \widetilde\omega_c^{n+1}<0}}
  \int_{S^2}c\sigma
  &={}\inf_{\substack{
    \widetilde\omega\in\mathcal T_0(\widetilde M^{n+1-k}),\ c\in\Omega^0(S^2)\\
    \widetilde\omega_c^{n+1}>0}}
  \int_{S^2}c\sigma\\
  &\ge\frac{\pi k}{n+1}.
\end{aligned}
\]
合并这两个不等式，得到所需的
$\varepsilon(\psi^k)\ge\pi$。定理得证。
\end{proof}

\begin{remark}\label{rem:9-6-6}
\begin{enumerate}[label=\textup{(\roman*)}]
\item \BookRecordIndex{T0028}Hofer 范数在 $\operatorname{Ham}(M)$ 上给出一个\BookRecordIndex{T0313}双不变 Finsler 度量。（人们常稍微
滥用语言，把这个度量中从恒等元到
$\phi\in\operatorname{Ham}(M)$ 的距离
$d(\operatorname{id},\phi)$ 称为 $\phi$ 的\emph{\BookRecordIndex{T0028}Hofer 范数}。）在 Ostrover--Wagner
\srcite{317} 工作的基础上，Buhovsky 与 Ostrover \srcite{51} 最近证明，它是唯一一个
$C^\infty$-连续的这种度量。不过，$\operatorname{Ham}(M,\omega)$ 上还有并非 Finsler 型
的其他双不变度量；比较习题~12.4.8 与 Viterbo \srcite{404}。

\item Chekanov 在 \srcite{61} 中把 \BookRecordIndex{T0027}Hofer 度量推广到\BookRecordIndex{T0219}拉格朗日子流形空间。他证明：若
$(P,\omega)$ 是闭\BookRecordIndex{T0358}辛流形，而 $\mathcal L(P,\omega,L)$ 是给定\BookRecordIndex{T0219}拉格朗日子流形 $L$ 在
哈密顿群 $\operatorname{Ham}(P,\omega)$ 作用下的轨道，则可把两个\BookRecordIndex{T0219}拉格朗日子流形
$L_0,L_1$ 之间的距离取为所有如下道路
$\{\phi_t\}_{t\in[0,1]}$ 的长度下确界：它们位于
$\operatorname{Ham}(P,\omega)$ 中，并满足
$\phi_0=\operatorname{id}$ 与 $\phi_1(L_0)=L_1$。

\item Polterovich 估计哈密顿环路长度的方法也能推广到拉格朗日情形。Akveld 与 Salamon
在 \srcite{11} 中考察带拉格朗日边界条件、以圆盘为底的辛\BookRecordIndex{T0279}纤维化，由此为\BookRecordIndex{T0219}拉格朗日子流形
的环路定义了类似不变量。在有利情形，相对\BookRecordIndex{T0143}格罗莫夫--威滕不变量给出拉格朗日环路的一个
给定\BookRecordIndex{T0158}哈密顿同痕类中 Hofer 长度的下界。因此，对 $\mathbb{CP}^n$ 中实射影空间的环路，可以
仿照定理~\ref{thm:9-6-5} 得到显式长度估计。

\item 用 $(M,\omega)$ 的辛微分同胚之图代替该辛微分同胚，可把哈密顿群嵌入
$(M\times M,\pi_2^*\omega-\pi_1^*\omega)$ 的嵌入\BookRecordIndex{T0219}拉格朗日子流形空间。Ostrover
\srcite{315} 最近给出的一个优美例子说明，包含
\[
  \operatorname{Ham}(M,\omega)
  \lhook\joinrel\longrightarrow
  \mathcal L(M\times M,\pi_2^*\omega-\pi_1^*\omega,\operatorname{diag})
\]
会严重扭曲度量：他构造了
$\operatorname{Ham}(\mathbb T^{2n},\omega_0)$ 中一个直径无限的集合，而它在拉格朗日
空间中的像却有界。尽管如此，Akveld--Salamon \srcite{11} 的结果表明，Polterovich 的
环路长度估计足够稳健，至少在某些情形进入拉格朗日语境后仍然有效。Usher
\srcite{400} 描述了其他一些情形，其中 $\operatorname{Ham}(M)$ 的子集嵌入拉格朗日
空间时，度量不发生扭曲。

\item 命题~\ref{prop:9-6-4} 的证明中的论证已由 McDuff--Tolman \srcite{279} 推广到
许多其他圆周作用；例~11.4.8 与例~11.4.11 提到了其中一些推论。这些结果表明，如下意义
下的\BookRecordIndex{T0059}半自由哈密顿圆周作用在所有同伦环路中使 Hofer 长度最小：每一点的稳定子要么是
平凡群，要么是整个圆周。类似地，对充分小的 $T>0$，哈密顿群中由与时间无关的哈密顿
函数生成的道路
$\{\phi_t^H\}_{0\le t\le T}$ 是绝对长度最短的；见 McDuff \srcite{271}。这些证明把
\BookRecordIndex{T0028}Hofer 范数与第~\ref{sec:9-1} 节\BookRecordIndex{T0419}作用谱
$\operatorname{Spec}(H)$ 的某些不变量等同起来。
\footnote{一般而言，\BookRecordIndex{T0028}Hofer 范数与\BookRecordIndex{T0272}谱范数不同；但在这个特殊情形二者相同。}
$\operatorname{Spec}(H)$ 的性质最好像第~12.4 节那样在 弗洛尔理论中研究。

\item \BookRecordIndex{T0161}哈密顿辛微分同胚群 $\operatorname{Ham}(M,\omega)$ 的拓扑、几何和一般性质是一个
非常活跃的研究主题。一个有趣的公开问题是：该群相对于 \BookRecordIndex{T0027}Hofer 度量是否总有无限直径。
关于这个问题的近期结果，见 Usher \srcite{400} 与 McDuff \srcite{274}。另一组问题
涉及 $\operatorname{Ham}(M,\omega)$ 上\BookRecordIndex{T0253}拟同态的存在性与性质。其中一些（如
\srcite{105} 中的\BookRecordIndex{T0253}拟同态）由第~12.4 节讨论的\BookRecordIndex{T0271}谱不变量构造，另一些则来自其他微分几何
构造；例如见 Py \srcite{332} 与 Shelukhin \srcite{374}。Polterovich 及其合作者近期的
工作 \srcite{107,52,331} 表明，\BookRecordIndex{T0253}拟同态与量子物理有非常有趣的联系，并带来新的辛刚性
现象。Müller 与 Oh \srcite{308} 的工作催生了研究 $C^0$-辛拓扑的新方法，Viterbo
\srcite{407} 与 Buhovsky--Seyfaddini \srcite{53} 展示了这些方法。更多问题与参考文献见
Polterovich \srcite{330} 或 McDuff \srcite{272}。
\end{enumerate}
\end{remark}
